\documentclass[aspectratio=169,11pt]{beamer}

\usepackage[utf8]{inputenc}
\usepackage[T1]{fontenc}
\usepackage{textcomp}
\usepackage[french]{babel}
\usepackage{lmodern}
\usepackage{booktabs}
\usepackage{tabularx}
\usepackage{tikz}
\usetikzlibrary{positioning,arrows.meta}

\definecolor{mainblue}{RGB}{28,74,102}
\definecolor{darkgray}{RGB}{62,62,66}
\definecolor{ent}{RGB}{221,236,247}
\definecolor{tab}{RGB}{215,237,221}
\definecolor{mod}{RGB}{233,233,238}
\definecolor{ext}{RGB}{250,233,207}
\definecolor{alerte}{RGB}{250,226,226}
\definecolor{kored}{RGB}{176,42,42}
\definecolor{okgreen}{RGB}{46,125,50}

\setbeamertemplate{navigation symbols}{}
\setbeamercolor{frametitle}{fg=mainblue}
\setbeamercolor{title}{fg=mainblue}
\setbeamercolor{structure}{fg=mainblue}
\setbeamercolor{normal text}{fg=darkgray}
\setbeamerfont{frametitle}{series=\bfseries,size=\Large}
\setbeamertemplate{itemize items}[circle]
\setbeamertemplate{footline}{\hfill\scriptsize\color{darkgray}\insertframenumber\hspace{0.4cm}\vspace{0.25cm}}

\newcommand{\eur}{\,\texteuro}
\newcommand{\qui}[1]{\hfill{\scriptsize\color{darkgray}\mdseries #1}}
\newcommand{\amesurer}{{\color{kored}\textit{à mesurer}}}
\newcommand{\f}[1]{{\ttfamily\small #1}}
\newcolumntype{Y}{>{\raggedright\arraybackslash}X}

\tikzset{
  etage/.style={rectangle,draw=mainblue,fill=ent,rounded corners=2pt,align=center,
                font=\small,inner sep=4pt,minimum height=0.75cm,text width=4.6cm},
  cote/.style={rectangle,draw=darkgray!60,fill=mod,rounded corners=2pt,align=left,
               font=\scriptsize,inner sep=3pt,text width=6.3cm},
  acteur/.style={rectangle,draw=darkgray,fill=ext,rounded corners=2pt,align=center,
                 font=\small,inner sep=4pt,minimum height=0.65cm,text width=4.6cm},
  fl/.style={-{Latex[length=1.8mm]},semithick,draw=mainblue!80},
  lien/.style={densely dashed,draw=darkgray!60},
}

\title{Un assistant bancaire frugal, explicable, et qui agit}
\subtitle{Comprendre, personnaliser, conseiller --- avec le plus petit calcul qui suffit}
\author{Yahya \and Sohayb \and Younes}
\institute{Hackathon IA --- La Banque Postale $\times$ CGI $\times$ ENSEIRB-MATMECA --- Sujet 1}
\date{Mercredi 28 octobre 2026}

\begin{document}

\begin{frame}
  \titlepage
\end{frame}

% ─────────────────────────────── Yahya ───────────────────────────────
\begin{frame}{Le problème \qui{Yahya}}
  \begin{columns}[T]
    \column{0.52\textwidth}
    \begin{block}{Un client écrit à sa banque}
      \og Je veux acheter un appartement. \fg\\
      \og J'ai perdu ma carte hier soir, je fais quoi ? \fg\\
      \og Pourquoi j'ai payé 8\eur{} de frais ce mois-ci ? \fg
    \end{block}
    \medskip
    L'assistant doit \textbf{comprendre}, \textbf{poser les bonnes questions}, tenir compte de
    \textbf{la situation du client} et proposer une piste.
    \column{0.44\textwidth}
    \begin{itemize}
      \item Un Français sur quatre est client de La Banque Postale.
      \item L'expérience client est sa \textbf{priorité IA n\textsuperscript{o}\,1}.
      \item La règle : \textbf{open source, local}, le plus petit modèle qui suffit,
            frugalité \textbf{mesurée}.
    \end{itemize}
  \end{columns}
\end{frame}

\begin{frame}{Notre démarche, en quatre temps \qui{Yahya}}
  \begin{tabularx}{\textwidth}{@{}l Y@{}}
    \toprule
    \textbf{Cadrer} & Les 9 situations du sujet. Constat : Banking77 couvre le \emph{service}
      (carte, frais, virements), pas le \emph{conseil} (épargne, crédit, logement). \\
    \addlinespace
    \textbf{Préparer} & \textbf{Banking77 n'est pas traduit} : un encodeur multilingue comprend le
      français. 8 intentions créées ; un test français écrit à la main ; la colonne \f{duration}
      de Bank Marketing exclue (fuite). \\
    \addlinespace
    \textbf{Concevoir} & Une cascade : chaque étage ne passe la main que s'il ne sait pas faire. \\
    \addlinespace
    \textbf{Démontrer} & Tests officiels, scénarios de bout en bout, comparaison avec une version
      \og tout LLM \fg, démo en direct. \\
    \bottomrule
  \end{tabularx}
\end{frame}

\begin{frame}{L'architecture : une cascade \qui{Yahya}}
  \centering
  \resizebox{!}{0.8\textheight}{%
  \begin{tikzpicture}[node distance=0.22cm]
    \node[acteur] (msg) {Message du client};
    \node[etage,below=of msg] (e1) {\textbf{1. Comprendre}};
    \node[etage,below=of e1] (e2) {\textbf{2. Converser}};
    \node[etage,below=of e2] (e3) {\textbf{3. Conseiller} \textbar{} \textbf{4. Agir}};
    \node[etage,below=of e3] (e5) {\textbf{5. Formuler}};
    \node[acteur,below=of e5] (rep) {Réponse\\ + \og Pourquoi \fg{} + \og Frugalité \fg};
    \node[cote,right=0.5cm of e1] (c1) {réponse attendue $\rightarrow$ mots-clés $\rightarrow$ classifieur e5-small (118 M)};
    \node[cote,right=0.5cm of e2] (c2) {parcours YAML ; le profil n'est jamais redemandé};
    \node[cote,right=0.5cm of e3] (c3) {règles publiques (HCSF, plafonds) ; actions après un \og oui \fg};
    \node[cote,right=0.5cm of e5] (c5) {gabarit, ou \textbf{un seul} appel au petit LLM, sous garde-fou};
    \foreach \a/\b in {msg/e1,e1/e2,e2/e3,e3/e5,e5/rep} \draw[fl] (\a) -- (\b);
    \foreach \i in {1,2,3,5} \draw[lien] (e\i) -- (c\i);
  \end{tikzpicture}}
\end{frame}

\begin{frame}{Le moins cher d'abord \qui{Yahya}}
  \begin{columns}[T]
    \column{0.5\textwidth}
    \begin{block}{Comprendre sans LLM}
      \begin{enumerate}
        \item la réponse à la question en cours ;
        \item les mots-clés, s'ils sont sûrs ;
        \item le classifieur, au-dessus du seuil $\tau$ ;
        \item sinon : une question au client, puis un conseiller.
      \end{enumerate}
      Les messages incompris sont journalisés pour améliorer le système.
    \end{block}
    \column{0.46\textwidth}
    \begin{block}{Décider sans LLM}
      Des règles que le client peut lire ; chaque conseil dit \textbf{pourquoi}, et
      \textbf{pourquoi pas}.
    \end{block}
    \begin{block}{Le LLM : une fois, à la fin}
      Il met en mots des faits calculés. Le garde-fou refuse tout chiffre inventé ou oublié.
    \end{block}
  \end{columns}
  \medskip
  \centering\small\textit{\og Le vélo pour aller chercher le pain, la voiture pour le week-end. \fg}
\end{frame}

\begin{frame}{Démonstration en direct \qui{Yahya pilote, Sohayb commente}}
  \footnotesize
  \begin{tabularx}{\textwidth}{@{}l l Y@{}}
    \toprule
    \textbf{Client} & \textbf{Message} & \textbf{Ce qu'on montre} \\
    \midrule
    Inès, 34 ans & \og J'ai perdu ma carte \fg & Opposition après un \og oui \fg ; aucun LLM \\
    Léa, 31 ans & \og Je veux acheter un appartement \fg & Revenus non redemandés ; 41,3\,\% $>$ 35\,\% ;
      \og et avec 200\,000 ? \fg{} : 32,2\,\% \\
    Sophie / Lucas & \og Je voudrais mettre de côté \fg & Même phrase, deux clients, deux conseils \\
    Marc, 52 ans & \og Pourquoi 8\eur{} de frais ? \fg & Historique, plafond légal, offre clientèle fragile \\
    Julie & \og Votre appli ne marche jamais ! \fg & Excuses, passage de main avec résumé \\
    Paul, 68 ans & \og Envoyer 500\eur{} à mon fils \fg & Mots simples ; \og oui, mais \fg{} n'est pas un accord ; montant corrigé \\
    \bottomrule
  \end{tabularx}
  \smallskip
  \scriptsize Démo de secours : 17 tours rejoués par le vrai pipeline, 0 appel au LLM, moins de 3~ms par tour
  (gabarits, sur un portable).
\end{frame}

\begin{frame}{Choisir le LLM en descendant en taille \qui{Yahya}}
  \small
  Seuil fixé avant de mesurer : \textbf{zéro chiffre inventé ni oublié} et \textbf{27 réponses
  sur 30} jugées bonnes. On garde \textbf{le plus petit} qui tient le seuil.
  \medskip

  \centering\small
\begin{tabular}{@{}l r c c r r@{}}
\toprule
\textbf{Modèle} & \textbf{Taille} & \textbf{Garde-fou} & \textbf{Juges} & \textbf{Latence p50} & \textbf{Tokens} \\
\midrule
\texttt{ministral-3:14b} & 9,1 Go & 23/30 & \textit{à juger} & 2,6 s & 483 \\
\texttt{ministral-3:8b} & 6,0 Go & 27/30 & \textit{à juger} & 1,9 s & 478 \\
\texttt{gemma4:e2b} & 4,6 Go & 27/30 & \textit{à juger} & 0,6 s & 449 \\
\texttt{ministral-3:3b} & 3,0 Go & 22/30 & \textit{à juger} & 0,8 s & 461 \\
\texttt{qwen3:1.7b} & 1,4 Go & 30/30 & \textit{à juger} & 0,6 s & 484 \\
\texttt{gemma3:1b} & 0,8 Go & 25/30 & \textit{à juger} & 0,4 s & 419 \\
\bottomrule
\end{tabular}

\medskip
{\scriptsize Mesuré sur DeepEirb, RTX 6000 Ada 48 Go. Garde-fou : aucun chiffre inventé ni oublié, longueur plausible.}

\end{frame}

% ─────────────────────────────── Younes ───────────────────────────────
\begin{frame}{Résultats : comprendre \qui{Younes}}
  \scriptsize
\begin{tabularx}{\textwidth}{@{}Y c c c c@{}}
\toprule
\textbf{Configuration} & \textbf{Anglais (intention, 77)} & \textbf{Français (routage)} & \textbf{Latence p50} & \textbf{Taille} \\
\midrule
\textit{Référence : BERT affiné (article)} & \textit{93,7\,\%} & --- & --- & \textit{0,44 Go} \\
Mots-clés & --- & 21,2\,\% acceptés, 88,7\,\% justes & $<$ 0,1 ms & 1 Ko \\
TF-IDF + rég. logistique & 89,2\,\% & 66,8\,\% & 0,4 ms & 3,9 Mo \\
\textbf{e5-small + rég. logistique} & \textbf{92,3\,\%} & \textbf{84,8\,\%} & 140,9 ms & 470,8 Mo \\
\textbf{Étage 1} : mots-clés puis e5 ($\tau$) & --- & 76,4\,\% compris, 93,7\,\% justes & 138,7 ms & --- \\
LLM 14B, sans exemple & \amesurer & \amesurer & \amesurer & 9,1 Go \\
\bottomrule
\end{tabularx}

\smallskip
{\tiny Test anglais : 3\,080 phrases Banking77 ; test français : 250 phrases écrites par l'équipe, jamais vues à l'entraînement. Routage : bonne famille pour une demande de service, bonne intention pour nos parcours. Mesuré sur yahya-bel-hajjam (x86\_64).}

  \medskip

  \small Entraîné sur Banking77 \textbf{en anglais, sans traduction} : l'encodeur multilingue
  transfère au français. Sous $\tau$, une question fermée ; jamais de devinette.
\end{frame}

\begin{frame}{Résultats : frugalité \qui{Younes}}
  \scriptsize
\begin{tabularx}{\textwidth}{@{}Y c c c c c@{}}
\toprule
\textbf{Configuration} & \textbf{Scénarios} & \textbf{Wh / 1\,000 req.} & \textbf{Latence p95} & \textbf{Go} & \textbf{Sans LLM} \\
\midrule
A --- un LLM fait tout (\texttt{ministral-3:14b}) & à juger & 115,41 & 2299 ms & 9,1 & 0\,\% \\
\textbf{B --- notre aiguilleur (niveaux 0, 1, 2)} (\texttt{qwen3:1.7b}) & 15 / 15 & 5,53 & 527 ms & 1,9 & 76\,\% \\
C --- niveau 0 seul (règles, sans LLM) & 15 / 15 & 0,178 & 19 ms & 0,5 & 100\,\% \\
\bottomrule
\end{tabularx}

\smallskip
{\tiny Énergie : CodeCarbon hors ligne, lots de 200 requêtes × 3 (A : 40 × 1) après chauffe, puissance de repos soustraite. Mesuré sur DeepEirb, RTX 6000 Ada 48 Go. Aiguilleur : niveau 0 : 76\,\%, niveau 1 : 12\,\%, niveau 2 : 12\,\% des messages.}

  \medskip

  \small La cascade répond sans LLM dans la plupart des tours ; le LLM ne reformule que les conseils.
\end{frame}

% ─────────────────────────────── Sohayb ───────────────────────────────
\begin{frame}{Explicabilité : pourquoi, et pourquoi pas \qui{Sohayb}}
  \begin{columns}[T]
    \column{0.5\textwidth}
    \textbf{Léa} : 250\,000\eur{}, 25\,000\eur{} d'apport, 2\,900\eur{} nets.
    \begin{itemize}
      \item Mensualité $\approx$ 1\,197\eur{} sur 25 ans à 3,32\,\% (taux moyen Banque de France,
            daté), hors assurance.
      \item \textcolor{kored}{\textbf{41,3\,\% $>$ 35\,\%}} : règle \f{HCSF-taux-effort} ; non retenu,
            sans refus bancaire.
      \item Écarté : prêt sur 30 ans (\f{HCSF-duree}).
      \item Budget indicatif $\approx$ 215\,000\eur{} ; à 200\,000\eur{}, 32,2\,\% : un conseiller étudie.
    \end{itemize}
    \column{0.46\textwidth}
    \begin{block}{Ce que voit le client}
      Les règles appliquées, les produits proposés \emph{et écartés}, la source de chaque
      paramètre. Le calcul se vérifie à la main.
    \end{block}
  \end{columns}
\end{frame}

\begin{frame}{Confiance et conformité \qui{Sohayb}}
  \begin{itemize}
    \item \textbf{Règlement européen sur l'IA} : la notation de solvabilité est à haut risque ;
          nous n'en faisons pas, un conseiller décide. Le client sait qu'il parle à une IA.
    \item \textbf{Mention légale} sur toute réponse de crédit ; aucune promesse.
    \item \textbf{Aucune action sans un \og oui \fg} ; rappel : jamais de code confidentiel.
    \item \textbf{Données simulées, traitement local} : rien ne sort de la machine.
    \item \textbf{Inclusion} : mots simples, jargon expliqué, offre clientèle fragile.
  \end{itemize}
\end{frame}

\begin{frame}{Limites et suite \qui{Younes}}
  \begin{columns}[T]
    \column{0.48\textwidth}
    \begin{block}{Ce qui est simulé}
      Profils et historiques (Bank Marketing enrichi), système bancaire, taux indicatifs.
    \end{block}
    \column{0.48\textwidth}
    \begin{block}{Pour 17\,000 bureaux}
      Coût par requête mesuré, donc extrapolable ; compréhension et règles sur CPU ; un serveur
      partagé pour le petit LLM ; voix locale et réponse dans la langue du client.
    \end{block}
  \end{columns}
\end{frame}

\begin{frame}
  \centering
  \vfill
  {\Large\color{mainblue}\textbf{Merci}}\par
  \bigskip
  \begin{minipage}{0.82\textwidth}\centering
    \textit{Notre assistant comprend avec de simples mots-clés ou un modèle de 118 millions de
    paramètres, décide avec des règles que le client peut lire, et n'appelle un LLM qu'une fois,
    à la fin, pour mettre en mots ce qui est déjà calculé.}
  \end{minipage}
  \vfill
\end{frame}

\end{document}
